
\begin{figure*}[t]
    \centering
     \includegraphics[width=\textwidth]{figures/method.pdf}
    \caption{We observe that we can find 1\% of the salient weights in LLMs based on the \emph{activation distribution} (middle). Keeping the salient weights in FP16 can significantly improve the quantized performance (PPL from 43.2 (left) to 13.0 (middle)), but the mixed-precision format is not hardware-efficient. We follow the activation-awareness principle and propose \methodshort (right). \methodshort performs per-channel scaling to protect the salient weights and reduce quantization error. We measure the perplexity of OPT-6.7B under INT3-g128 quantization. %
    } 
    \label{fig:overview}
\end{figure*}


\section{Related Work}
\myparagraph{Model quantization methods. }
Quantization reduces the bit-precision of deep learning models~\cite{han2016deep, jacob2018quantization, nagel2019data, wang2019haq, nagel2020up,lin2020mcunet}, which helps to reduce the model size and accelerate inference. Quantization techniques generally fall into two categories: quantization-aware training (QAT, which relies on backpropagation to update the quantized weights)~\cite{bengio2013estimating, gholami2021survey, nagel2021white, choi2018pact} and post-training quantization~\cite{jacob2018quantization, nagel2019data, nagel2020up} (PTQ, usually training-free). The QAT methods cannot easily scale up to large models like LLMs. Therefore, people usually use PTQ methods to quantize LLMs. 

\myparagraph{Quantization of LLMs.} People study two settings for LLM quantization: (1) W8A8 quantization, where both activation and weights are quantized to INT8~\cite{dettmers2022llmint8, xiao2022smoothquant, zeroquant, outlier_suppression, wei2023outlier}; (2) Low-bit weight-only quantization (\eg, W4A16), where only weights are quantized into low-bit integers~\cite{frantar2022gptq, dettmers2022case, sheng2023high, park2022nuqmm}. We focus on the second setting in this work since it not only reduces the hardware barrier (requiring a smaller memory size) but also speeds up the token generation (remedies memory-bound workload). Apart from the vanilla round-to-nearest baseline (RTN), GPTQ~\cite{frantar2022gptq} is the closest to our work. However, the reconstruction process of GPTQ leads to an over-fitting issue to the calibration set and may not preserve the generalist abilities of LLMs for other modalities and domains. It also requires a reordering trick to work for some models (\eg, LLaMA-7B~\cite{touvron2023llama} and OPT-66B~\cite{opt}). Apart from quantiztion methods designed for general-purporse hardware, SpAtten~\cite{spatten} designs a progressive approach to gradually increase the number of bits used in softmax calculation.

\myparagraph{System support for low-bit quantized LLMs. } Low-bit quantized LLMs have been a popular setting to reduce inference costs. There are some system supports to achieve a practical speed-up. GPTQ~\cite{frantar2022gptq} provides INT3 kernels for OPT models and \texttt{GPTQ-for-LLaMA} extends kernel support for INT4 reordered quantization with the help of Triton~\cite{tillet2019triton}. FlexGen~\cite{sheng2023high}, \texttt{llama.cpp}\footnote{https://github.com/ggerganov/llama.cpp} and \texttt{exllama}\footnote{https://github.com/turboderp/exllama} perform group-wise INT4 quantization to reduce I/O costs and offloading. FasterTransformer implements FP16$\times$INT4 GEMM for weight-only per-tensor quantization but does not support group quantization. LUT-GEMM~\cite{park2022nuqmm} performs bitwise computation on GPU CUDA cores with the help of lookup tables. Our concurrent work, MLC-LLM~\cite{mlc-llm} offers strong results on multiple edge CPU and GPU platforms thanks to the powerful TVM~\cite{chen2018tvm,feng2022tensorir} backend. %


